% GENERATED FILE. Edit canonical lessons, then run npm run generate:books.
% canonical-source-hash: fnv1a64:d3081ca24d7e55e8
% canonical-lessons: JA-C187-azukeru, JA-C187-orosu, JA-C187-tsuuchou, JA-C187-kariru, JA-C187-ryougae

\chapter{At the Bank}
\label{ch:ja-at-the-bank}
\hlchaptermodality{187}

\begin{chapteropening}
\textbf{By the end of this chapter:} \emph{I can deposit, withdraw and change money at the bank.}

Complete the last lesson of chapter 187: I can deposit, withdraw and change money at the bank.
\end{chapteropening}

\section[\texorpdfstring{\ja{あずける} (azukeru)}{あずける (azukeru)}]{\ja{あずける} (azukeru) — to deposit; to leave in someone's care}

\label{lesson:JA-C187-azukeru}



\begin{quote}
Before the new one: say the Japanese for fee, then the Japanese for tax.
\end{quote}

\subsection*{You'll want to know: \ja{あずける}}
\textbf{\ja{あずける}} — \emph{azukeru} — \textquotedblleft{}to deposit; to leave in someone's care\textquotedblright{}.

\begin{quote}
\textbf{\ja{ぎんこうに} \ja{おかねを} \ja{あずけます。}} — \emph{ginkō ni okane o azukemasu} — \textquotedblleft{}I deposit money in the bank.\textquotedblright{}
\end{quote}

Polite \textbf{\ja{あずけます}}, \emph{azukemasu}; an ichidan verb.

\begin{cousinweb}
\textbf{\ja{あずける}} (\emph{azukeru}) is a native Japanese word (wago), inherited rather than borrowed from Chinese; its kanji was attached later as a kun'yomi (JMdict, KANJIDIC2).
\end{cousinweb}

\begin{grammarlens}[title={What you've built}]
One more word for everyday life at A2.
\end{grammarlens}

\subsection*{Your turn}
\begin{itemize}
\raggedright
  \item \emph{Say it:} \emph{azukeru}
  \item \emph{Say it:} \emph{azukeru}, once more
  \item \emph{Say it:} \emph{ginkō ni okane o azukemasu} — I deposit money in the bank.
  \item \emph{Recall:} say the Japanese for fee, then the Japanese for tax, then say \emph{azukeru} again
\end{itemize}

\begin{tcolorbox}[breakable,colback=teal!4,colframe=teal!35!black,title={Before you move on}]
What does \ja{あずける} mean? (\textquotedblleft{}To deposit; to leave in someone's care\textquotedblright{}.) Say it once more.
\end{tcolorbox}
\section[\texorpdfstring{\ja{おろす} (orosu)}{おろす (orosu)}]{\ja{おろす} (orosu) — to withdraw (money); to take down}

\label{lesson:JA-C187-orosu}



\begin{quote}
Before the new one: say the Japanese for tax, then the Japanese for to deposit.
\end{quote}

\subsection*{You'll want to know: \ja{おろす}}
\textbf{\ja{おろす}} — \emph{orosu} — \textquotedblleft{}to withdraw (money); to take down\textquotedblright{}.

\begin{quote}
\textbf{\ja{ぎんこうで} \ja{おかねを} \ja{おろします。}} — \emph{ginkō de okane o oroshimasu} — \textquotedblleft{}I withdraw money at the bank.\textquotedblright{}
\end{quote}

Polite \textbf{\ja{おろします}}, \emph{oroshimasu}; a godan verb.

\begin{cousinweb}
\textbf{\ja{おろす}} (\emph{orosu}) is a native Japanese word (wago), inherited rather than borrowed from Chinese; its kanji was attached later as a kun'yomi (JMdict, KANJIDIC2).
\end{cousinweb}

\begin{grammarlens}[title={What you've built}]
One more word for everyday life at A2.
\end{grammarlens}

\subsection*{Your turn}
\begin{itemize}
\raggedright
  \item \emph{Say it:} \emph{orosu}
  \item \emph{Say it:} \emph{orosu}, once more
  \item \emph{Say it:} \emph{ginkō de okane o oroshimasu} — I withdraw money at the bank.
  \item \emph{Recall:} say the Japanese for tax, then the Japanese for to deposit, then say \emph{orosu} again
\end{itemize}

\begin{tcolorbox}[breakable,colback=teal!4,colframe=teal!35!black,title={Before you move on}]
What does \ja{おろす} mean? (\textquotedblleft{}To withdraw (money); to take down\textquotedblright{}.) Say it once more.
\end{tcolorbox}
\section[\texorpdfstring{\ja{つうちょう} (tsūchō)}{つうちょう (tsūchō)}]{\ja{つうちょう} (tsūchō) — bankbook; passbook}

\label{lesson:JA-C187-tsuuchou}



\begin{quote}
Before the new one: say the Japanese for to deposit, then the Japanese for to withdraw.
\end{quote}

\subsection*{You'll want to know: \ja{つうちょう}}
\textbf{\ja{つうちょう}} — \emph{tsūchō} — \textquotedblleft{}bankbook; passbook\textquotedblright{}.

\begin{quote}
\textbf{\ja{ぎんこうで} \ja{つうちょうを} \ja{だします。}} — \emph{ginkō de tsūchō o dashimasu} — \textquotedblleft{}I hand in my bankbook at the bank.\textquotedblright{}
\end{quote}

\begin{cousinweb}
\textbf{\ja{つうちょう}} (\emph{tsūchō}) is Sino-Japanese (kango): two Chinese characters read with their on'yomi, readings borrowed from Chinese. They are \emph{tsū}, \textquotedblleft{}pass through\textquotedblright{} (Middle Chinese \emph{thuwng}) and \emph{chō}, \textquotedblleft{}notebook, account book\textquotedblright{} (\emph{trjangH}) (KANJIDIC2; Middle Chinese in Baxter's transcription).
\end{cousinweb}

\begin{grammarlens}[title={What you've built}]
One more word for everyday life at A2.
\end{grammarlens}

\subsection*{Your turn}
\begin{itemize}
\raggedright
  \item \emph{Say it:} \emph{tsūchō}
  \item \emph{Say it:} \emph{tsūchō}, once more
  \item \emph{Say it:} \emph{ginkō de tsūchō o dashimasu} — I hand in my bankbook at the bank.
  \item \emph{Recall:} say the Japanese for to deposit, then the Japanese for to withdraw, then say \emph{tsūchō} again
\end{itemize}

\begin{tcolorbox}[breakable,colback=teal!4,colframe=teal!35!black,title={Before you move on}]
What does \ja{つうちょう} mean? (\textquotedblleft{}Bankbook; passbook\textquotedblright{}.) Say it once more.
\end{tcolorbox}
\section[\texorpdfstring{\ja{かりる} (kariru)}{かりる (kariru)}]{\ja{かりる} (kariru) — to borrow; to rent}

\label{lesson:JA-C187-kariru}



\begin{quote}
Before the new one: say the Japanese for to withdraw, then the Japanese for bankbook.
\end{quote}

\subsection*{You'll want to know: \ja{かりる}}
\textbf{\ja{かりる}} — \emph{kariru} — \textquotedblleft{}to borrow; to rent\textquotedblright{}.

\begin{quote}
\textbf{\ja{あにに} \ja{おかねを} \ja{かります。}} — \emph{ani ni okane o karimasu} — \textquotedblleft{}I borrow money from my older brother.\textquotedblright{}
\end{quote}

Polite \textbf{\ja{かります}}, \emph{karimasu}; an ichidan verb. Its partner \textbf{\ja{かす}}, \textquotedblleft{}to lend\textquotedblright{}, is taught already.

\begin{cousinweb}
\textbf{\ja{かりる}} (\emph{kariru}) is a native Japanese word (wago), inherited rather than borrowed from Chinese; its kanji was attached later as a kun'yomi (JMdict, KANJIDIC2).
\end{cousinweb}

\begin{grammarlens}[title={What you've built}]
One more word for everyday life at A2.
\end{grammarlens}

\subsection*{Your turn}
\begin{itemize}
\raggedright
  \item \emph{Say it:} \emph{kariru}
  \item \emph{Say it:} \emph{kariru}, once more
  \item \emph{Say it:} \emph{ani ni okane o karimasu} — I borrow money from my older brother.
  \item \emph{Recall:} say the Japanese for to withdraw, then the Japanese for bankbook, then say \emph{kariru} again
\end{itemize}

\begin{tcolorbox}[breakable,colback=teal!4,colframe=teal!35!black,title={Before you move on}]
What does \ja{かりる} mean? (\textquotedblleft{}To borrow; to rent\textquotedblright{}.) Say it once more.
\end{tcolorbox}
\section[\texorpdfstring{\ja{りょうがえ} (ryōgae)}{りょうがえ (ryōgae)}]{\ja{りょうがえ} (ryōgae) — money exchange; changing money}

\label{lesson:JA-C187-ryougae}



\begin{quote}
Before the new one: say the Japanese for bankbook, then the Japanese for to borrow.
\end{quote}

\subsection*{You'll want to know: \ja{りょうがえ}}
\textbf{\ja{りょうがえ}} — \emph{ryōgae} — \textquotedblleft{}money exchange; changing money\textquotedblright{}.

\begin{quote}
\textbf{\ja{ぎんこうで} \ja{りょうがえを} \ja{します。}} — \emph{ginkō de ryōgae o shimasu} — \textquotedblleft{}I change money at the bank.\textquotedblright{}
\end{quote}

\begin{cousinweb}
\textbf{\ja{りょうがえ}} (\emph{ryōgae}) is built from Sino-Japanese \emph{ryō}, an old Japanese coin, a reading borrowed from Chinese (Middle Chinese \emph{ljangX}), and native \emph{kae}, \textquotedblleft{}exchange\textquotedblright{}, from \emph{kaeru}, \textquotedblleft{}to exchange\textquotedblright{} (JMdict, KANJIDIC2).
\end{cousinweb}

\begin{grammarlens}[title={What you've built}]
One more word for everyday life at A2.
\end{grammarlens}

\subsection*{Your turn}
\begin{itemize}
\raggedright
  \item \emph{Say it:} \emph{ryōgae}
  \item \emph{Say it:} \emph{ryōgae}, once more
  \item \emph{Say it:} \emph{ginkō de ryōgae o shimasu} — I change money at the bank.
  \item \emph{Recall:} say the Japanese for bankbook, then the Japanese for to borrow, then say \emph{ryōgae} again
\end{itemize}

\begin{tcolorbox}[breakable,colback=teal!4,colframe=teal!35!black,title={Before you move on}]
What does \ja{りょうがえ} mean? (\textquotedblleft{}Money exchange; changing money\textquotedblright{}.) Say it once more.
\end{tcolorbox}
